% ============ MỤC THỰC NGHIỆM (bản thay thế) ============
% Thay thế toàn bộ từ \section{Thực nghiệm} đến trước \section{Nhận xét và kết luận}
% trong report. Phần "Chỉ số đánh giá" giữ nguyên như bản cũ (đã đầy đủ) — chèn lại
% nguyên văn tại vị trí đánh dấu bên dưới.
% Các hình dùng đường dẫn img/report/*.png làm placeholder — thay bằng ảnh thật.

\section{Thực nghiệm}

Mục tiêu của phần thực nghiệm là huấn luyện lại và đánh giá bốn mô hình đã trình
bày ở trên (SRCNN, SRGAN, SwinIR, SR3) trong cùng một điều kiện được kiểm soát,
từ đó so sánh trực tiếp bốn hướng tiếp cận: CNN, GAN, Transformer và diffusion.
Toàn bộ thực nghiệm chạy trên một máy cá nhân với thông số như Bảng~\ref{tab:environment}.

\begin{table}[H]
\centering
\caption{Thông số môi trường thực nghiệm.}
\label{tab:environment}
\begin{tabular}{|l|l|}
\hline
\textbf{Thông số} & \textbf{Chi tiết} \\ \hline
OS & Linux 6.8.0 (x86\_64) \\ \hline
Python & 3.12 \\ \hline
Framework & PyTorch 2.12 (CUDA 13) \\ \hline
CPU & 16 cores (x86\_64) \\ \hline
RAM & 14 GB \\ \hline
GPU & NVIDIA GeForce RTX 4060 Laptop (8 GB, giới hạn dùng 6 GB) \\ \hline
Mixed precision & bfloat16 autocast (SwinIR, SR3) \\ \hline
\end{tabular}
\end{table}

Mã nguồn và hướng dẫn sử dụng được lưu giữ tại: \url{https://github.com/tannguyenbmt47/Image_Super_Resolution_evaluation}.

\subsection{Tập dữ liệu}
\label{sec:datasets}

\textbf{Dữ liệu huấn luyện.} Các mô hình được huấn luyện trên tập DF2K, là hợp của
DIV2K train (800 ảnh) và Flickr2K (2650 ảnh), tổng cộng 3450 ảnh high-resolution
chất lượng 2K với nội dung đa dạng. Ảnh low-resolution được sinh từ ảnh HR bằng
phép thu nhỏ bicubic có antialiasing với hệ số $\times 4$, theo đúng protocol
bicubic degradation chuẩn của bài toán classical SR. Để loại bỏ nút thắt giải mã
ảnh 2K trong quá trình nạp dữ liệu, toàn bộ ảnh HR được cắt trước thành 30581
tile kích thước $480 \times 480$ không chồng lấp; mỗi bước huấn luyện cắt ngẫu
nhiên một patch từ tile, kèm augmentation lật ngang/dọc và xoay $90^{\circ}$.

\textbf{Dữ liệu đánh giá.} Việc đánh giá thực hiện trên DIV2K validation (100
ảnh), tách riêng hoàn toàn khỏi dữ liệu huấn luyện để điểm số phản ánh khả năng
tổng quát hóa. Do chi phí self-attention của mô hình diffusion không cho phép xử
lý ảnh 2K nguyên bản, tất cả mô hình được đánh giá trên cùng vùng center-crop
$256 \times 256$ của mỗi ảnh nhằm bảo đảm so sánh thống nhất.

% ================================================================
% CHÈN LẠI NGUYÊN VĂN \subsection{Chỉ số đánh giá} (PSNR/SSIM/LPIPS/
% NIQE/MUSIQ/CLIPIQA + bảng tổng hợp) TỪ BẢN CŨ VÀO ĐÂY.
% ================================================================

\subsection{Thiết lập huấn luyện}
\label{sec:training-setup}

Bốn mô hình thuộc bốn họ kiến trúc khác nhau, vì vậy giao thức so sánh công bằng
được cố định như sau:

\begin{itemize}
\item \textbf{Ngân sách dữ liệu như nhau:} mỗi mô hình được huấn luyện 6 epoch
      trên 30581 tile, tương đương khoảng 50 epoch trên 3450 ảnh gốc
      ($\approx 183$ nghìn patch). Early stopping được kích hoạt khi chỉ số
      validation không cải thiện sau 3 lần đánh giá liên tiếp.
\item \textbf{Cùng dữ liệu và degradation:} cùng tập DF2K, cùng bicubic $\times 4$.
\item \textbf{Cùng giao thức đánh giá:} cùng DIV2K-Val center-crop 256, cùng sáu
      chỉ số.
\end{itemize}

Siêu tham số của từng mô hình bám theo bài báo gốc và được tóm tắt trong
Bảng~\ref{tab:hparams}. SRCNN dùng MSE loss trên ảnh LR đã phóng to sẵn bằng
bicubic (pre-upsampling). SRGAN dùng content loss VGG54 với trọng số
$6 \times 10^{-3}$ kết hợp adversarial loss $10^{-3}$ và pixel loss MSE; generator
và discriminator được tối ưu xen kẽ. SwinIR dùng phiên bản lightweight (embed dim
60, 4 khối RSTB, cửa sổ $8 \times 8$) với L1 loss; phiên bản M (11.9M tham số)
không khả thi để huấn luyện từ đầu trong ngân sách bộ nhớ 6 GB. SR3 huấn luyện
U-Net dự đoán nhiễu với L1 loss, 1000 bước diffusion khi huấn luyện và 100 bước
DDIM khi suy diễn. Tất cả dùng Adam; learning rate giảm một nửa tại giữa quá
trình huấn luyện. SwinIR và SR3 dùng mixed precision bfloat16 để tăng gấp đôi
batch size trong cùng giới hạn bộ nhớ.

\begin{table}[H]
\centering
\caption{Cấu hình huấn luyện của bốn mô hình.}
\label{tab:hparams}
\begin{tabular}{|l|c|c|c|c|c|}
\hline
\textbf{Mô hình} & \textbf{Tham số} & \textbf{Loss} & \textbf{Batch} & \textbf{Learning rate} & \textbf{Patch HR} \\ \hline
SRCNN        & 0.07M  & MSE                       & 16 & $10^{-4}$          & 128 \\ \hline
SRGAN        & 6.76M\textsuperscript{*} & MSE + VGG54 + adv. & 16 & $10^{-4}$ & 96  \\ \hline
SwinIR-light & 0.93M  & L1                        & 12 & $2 \times 10^{-4}$ & 256 \\ \hline
SR3          & 63.15M & L1 (noise prediction)     & 6  & $10^{-4}$          & 160 \\ \hline
\multicolumn{6}{|l|}{\footnotesize \textsuperscript{*}Gồm generator 1.55M và discriminator; suy diễn chỉ dùng generator.} \\ \hline
\end{tabular}
\end{table}

Thời gian huấn luyện thực tế: SRCNN 9 phút, SRGAN 29 phút, SwinIR 1 giờ 28 phút,
SR3 2 giờ 15 phút (dừng sớm tại epoch 4 do early stopping). Tổng thời gian cho cả
bốn mô hình khoảng 4 giờ 20 phút.

\subsection{Kết quả đạt được}
\label{sec:results}

Bảng~\ref{tab:experiment-results} trình bày kết quả đánh giá trên DIV2K-Val theo
sáu chỉ số, kèm baseline nội suy bicubic để đối chiếu. PSNR, SSIM, MUSIQ và
CLIPIQA càng cao càng tốt; LPIPS và NIQE càng thấp càng tốt.

\begin{table}[H]
\centering
\caption{Kết quả đánh giá các mô hình super-resolution trên DIV2K-Val
($\times 4$, center-crop 256). \textbf{Đậm}: tốt nhất theo từng chỉ số.}
\label{tab:experiment-results}
\begin{tabular}{|l|c|c|c|c|c|c|}
\hline
\textbf{Model} & \textbf{PSNR} & \textbf{SSIM} & \textbf{LPIPS} & \textbf{NIQE} & \textbf{MUSIQ} & \textbf{CLIPIQA} \\ \hline
Bicubic & 26.81 & 0.7227 & 0.4716 & 8.94 & 22.63 & 0.5280 \\ \hline
SRCNN   & 27.49 & 0.7523 & 0.3999 & 8.43 & 27.25 & \textbf{0.5841} \\ \hline
SRGAN   & 26.00 & 0.6999 & \textbf{0.3031} & 5.60 & \textbf{38.53} & 0.4317 \\ \hline
SwinIR  & \textbf{28.01} & \textbf{0.7708} & 0.3632 & 8.35 & 29.56 & 0.4682 \\ \hline
SR3     & 19.32 & 0.6556 & 0.3586 & \textbf{4.85} & 30.27 & 0.4015 \\ \hline
\end{tabular}
\end{table}

Kết quả thể hiện rõ sự đánh đổi giữa độ trung thực (fidelity) và chất lượng cảm
nhận (perceptual quality) đã phân tích ở các mục trước:

\begin{itemize}
\item \textbf{SwinIR} đạt PSNR và SSIM cao nhất (28.01 dB / 0.7708). Với chỉ
      0.93M tham số, phiên bản lightweight vượt SRCNN 0.52 dB, cho thấy cơ chế
      window attention sử dụng ngân sách huấn luyện hiệu quả nhất cho mục tiêu
      tối thiểu hóa sai số điểm ảnh.
\item \textbf{SRGAN} chấp nhận giảm PSNR (thấp hơn bicubic 0.8 dB) để đổi lấy
      chất lượng cảm nhận: LPIPS tốt nhất (0.3031) và MUSIQ tốt nhất (38.53).
      Đây là hành vi đặc trưng của adversarial training: thay vì trung bình hóa
      các khả năng như mô hình tối ưu MSE, generator sinh texture sắc nét khiến
      ảnh trông tự nhiên hơn dù sai số điểm ảnh lớn hơn.
\item \textbf{SR3} đạt NIQE tốt nhất (4.85), nghĩa là ảnh sinh ra có thống kê
      gần ảnh tự nhiên nhất, nhưng PSNR thấp nhất (19.32 dB). Nguyên nhân thứ
      nhất là bản chất sinh mẫu ngẫu nhiên của diffusion khiến chi tiết có thể
      lệch khỏi ground truth; nguyên nhân thứ hai là mô hình 63M tham số còn xa
      điểm hội tụ với ngân sách huấn luyện này --- nhất quán với hạn chế đã nêu
      ở mục lý thuyết rằng diffusion đòi hỏi chi phí huấn luyện lớn hơn nhiều
      bậc so với các họ mô hình còn lại.
\item \textbf{SRCNN} vượt bicubic trên mọi chỉ số, xác nhận vai trò baseline hợp
      lệ; khoảng cách với SwinIR phản ánh giới hạn của kiến trúc 3 lớp tích chập.
\end{itemize}

\subsection{Đường cong huấn luyện}

Hình~\ref{fig:loss-curves} và Hình~\ref{fig:val-curves} trình bày diễn tiến huấn
luyện. Do bốn mô hình dùng bốn hàm loss khác bản chất (pixel loss, composite GAN
loss, noise-prediction loss), các đường loss chỉ dùng để quan sát tính hội tụ của
từng mô hình; việc so sánh giữa các mô hình dựa trên validation PSNR/SSIM đo trên
cùng tập DIV2K-Val.

\begin{figure}[H]
\centering
% PLACEHOLDER: xuất từ notebook train.ipynb (ô Visualize), hình 4 panel loss
\includegraphics[width=0.9\textwidth]{img/report/training-loss.png}
\caption{Training loss của bốn mô hình (mỗi mô hình một panel do thang đo loss khác nhau).}
\label{fig:loss-curves}
\end{figure}

\begin{figure}[H]
\centering
% PLACEHOLDER: xuất từ notebook train.ipynb (ô Visualize), hình val PSNR/SSIM
\includegraphics[width=0.9\textwidth]{img/report/val-psnr-ssim.png}
\caption{Validation PSNR và SSIM trên DIV2K-Val theo tiến trình huấn luyện.}
\label{fig:val-curves}
\end{figure}

\subsection{Kết quả định tính}

Hình~\ref{fig:qualitative} so sánh trực quan kết quả của các mô hình.
SwinIR khôi phục cạnh và cấu trúc ổn định nhất; SRGAN cho texture sắc nét hơn ở
vùng chi tiết cao (cỏ, tóc, vân gạch) nhưng đôi khi sinh chi tiết không có trong
ảnh gốc; SR3 cho ảnh nhìn tự nhiên nhưng nội dung chi tiết có thể sai lệch;
SRCNN mượt và an toàn nhưng thiếu chi tiết tần số cao.

\begin{figure}[H]
\centering
% PLACEHOLDER: lưới so sánh, mỗi hàng một ảnh, cột = LR / Bicubic / SRCNN / SRGAN / SwinIR / SR3 / Ground truth
\includegraphics[width=0.9\textwidth]{img/report/qualitative-grid.png}
\caption{So sánh định tính trên DIV2K-Val ($\times 4$): từ trái sang phải là ảnh
LR, Bicubic, SRCNN, SRGAN, SwinIR, SR3 và ground truth.}
\label{fig:qualitative}
\end{figure}

\subsection{Hạn chế của thực nghiệm}

Một số hạn chế cần lưu ý khi diễn giải kết quả: (1) ngân sách huấn luyện nhỏ hơn
nhiều so với các bài báo gốc (SwinIR gốc huấn luyện 500 nghìn iteration, SR3 gốc
khoảng 1 triệu bước), đặc biệt bất lợi cho SR3; (2) SwinIR dùng phiên bản
lightweight 0.93M thay vì phiên bản M 11.9M do giới hạn bộ nhớ GPU; (3) đánh giá
trên center-crop 256 thay vì ảnh đầy đủ do ràng buộc bộ nhớ của mô hình
diffusion; (4) thực nghiệm chỉ xét degradation bicubic, chưa xét các degradation
thực tế (blur, nhiễu, nén JPEG) và chưa thực hiện giai đoạn fine-tune trên dữ
liệu real-world như RealSR.
